\section{Introducción}

El análisis automático de manchas solares constituye un punto de partida en la física solar, ya que sienta las bases para el monitoreo operativo de la actividad fotosférica y para la prevención de eventos capaces de afectar sistemas electrónicos y de comunicación. Las regiones activas de la fotosfera concentran intensos campos magnéticos que exhiben estructuras oscuras en observaciones de luz blanca de observatorios terrestres y en órbita. En particular, en observatorios terrestres, la detección de grupos y la clasificación morfológica suelen ejecutarse como procedimientos independientes, con una elevada dependencia de la inspección manual; por ello, la trazabilidad entre la observación original y los datos de clasificación resulta limitada.

Por ello, este trabajo presenta la cadena operativa completa de adquisición, procesamiento y análisis predictivo de manchas solares. El prototipo de adquisición de imágenes de luz blanca emplea una cámara monocromática y un sistema de atenuación óptica con cuña de Herschel montados en un telescopio refractor Vixen (con un diámetro $D = 60$ mm y una distancia focal $F = 420$ mm), utilizando una Raspberry Pi 5 como nodo embebido para el preprocesamiento \citep{zwoASI174Manual2025, luntWhiteLightWedges2026, raspberryPi5Docs2026}.

Las imágenes son procesadas mediante una arquitectura de Deep Learning en cascada de visión por computadora basada en YOLO26n (versión 26 del modelo YOLO nano de Ultralytics) para localizar regiones candidatas y un clasificador jerárquico basado en ConvNeXtV2a (versión 2 del modelo ConvNeXt atto desarrollado por Meta AI) para estimar los caracteres $Z$, $P$ y $C$ \citep{ultralyticsYOLO26, woo2023convnextv2}. La salida se complementa con una etapa de normalización bayesiana que restringe los resultados obtenidos al conjunto de combinaciones taxonómicamente válidas \citep{mcintosh1990classification}.

Los registros generados enriquecen modelos estimadores del número de manchas solares (Sunspot Number o SSN) y de la distribución latitudinal de la actividad solar, implementados mediante un modelo PVL con filtro de Kalman adaptativo (PVLK) y un modelo espacio-temporal basado en la Descomposición de Modos Dinámicos (DMD) aplicada a la vectorización del diagrama de Spörer.

Los algoritmos desarrollados están integrados en una plataforma web que funciona como una interfaz unificada para el registro, la clasificación asistida y el análisis predictivo de los datos observados.